% called by main.tex
%
\chapter{Configuración experimental y reproducibilidad}
\label{ch::AnexoA}

Este anexo recopila los parámetros técnicos necesarios para reproducir los experimentos descritos en la memoria. La información se organiza en cuatro bloques: el entorno de
hardware y software, las características del modelo base y su cuantización, la configuración de LoRA compartida entre todas las variantes de ajuste fino, y los hiperparámetros
específicos de cada experimento. La motivación de las decisiones más relevantes aparece comentada junto a los valores correspondientes, de modo que el anexo sirva tanto como
referencia técnica como como justificación de las elecciones de implementación.

\section{Entorno de hardware y software}

Todos los experimentos de entrenamiento, validación e inferencia se realizaron en la misma máquina, sin uso de cómputo en la nube. Esta restricción determinó varias decisiones
de diseño: el uso de cuantización de cuatro bits para reducir la huella de memoria del modelo, la elección de un modelo de siete mil millones de parámetros en lugar de variantes
mayores, y el tamaño efectivo de batch. El hardware utilizado se detalla en la tabla~\ref{tab:AnexoA_hardware}.

\begin{table}[htbp]
\centering
\caption{Especificaciones del entorno de hardware.}
\label{tab:AnexoA_hardware}
\small
\begin{tabularx}{\textwidth}{lX}
\toprule
\textbf{Componente} & \textbf{Especificación} \\
\midrule
CPU & Intel Core i7-12700H (12.ª gen.), 14 núcleos físicos, 20 hilos lógicos \\
RAM & 16 GB DDR4 a 3200 MHz (2 $\times$ 8 GB Samsung) \\
GPU & NVIDIA GeForce RTX 3060 Laptop GPU, 6 GB VRAM GDDR6 \\
VRAM en uso (inferencia) & $\approx$ 4.8 GB \\
Sistema operativo & Windows 11 \\
\bottomrule
\end{tabularx}
\end{table}

El entorno de software se gestiona con \texttt{uv} como gestor de dependencias, con un entorno virtual aislado por proyecto. La tabla~\ref{tab:AnexoA_software} recoge las versiones
mínimas requeridas en el fichero \texttt{pyproject.toml}. Las dependencias de aceleración de cómputo se instalan desde el índice oficial de PyTorch con soporte CUDA~12.1.

\begin{table}[htbp]
\centering
\caption{Versiones mínimas de las dependencias principales.}
\label{tab:AnexoA_software}
\begin{tabular}{ll}
\toprule
\textbf{Librería} & \textbf{Versión mínima} \\
\midrule
Python & $\geq$ 3.12 \\
PyTorch & $\geq$ 2.5.0 (CUDA 12.1) \\
Transformers (HuggingFace) & $\geq$ 4.56.1 \\
PEFT & $\geq$ 0.17.1 \\
bitsandbytes & $\geq$ 0.43.0 \\
TRL & $\geq$ 0.23.1 \\
accelerate & $\geq$ 1.10.1 \\
wandb & $\geq$ 0.19.0 \\
\bottomrule
\end{tabular}
\end{table}

\section{Modelo base y cuantización}

El modelo base utilizado en todos los experimentos es \textbf{Qwen2.5-7B-Instruct}, en su versión cuantizada a cuatro bits. La elección de este modelo responde a tres criterios.
Primero, el tamaño de siete mil millones de parámetros permite su carga y ajuste en una GPU de consumo con seis gigabytes de VRAM gracias a la cuantización. Segundo, la variante
\textit{instruct} ha sido ajustada previamente con instrucciones diversas en varios idiomas y dominios técnicos, lo que proporciona un punto de partida con conocimiento razonable 
del vocabulario de Kubernetes sin necesidad de preentrenamiento adicional. Tercero, la familia Qwen2.5 ofrece una arquitectura \textit{decoder-only} estándar, que facilita la
integración de cabezales adicionales sobre los estados ocultos del transformer sin modificaciones profundas al flujo de inferencia.

La arquitectura interna del modelo se describe en la tabla~\ref{tab:AnexoA_model}. La cuantización se aplica mediante bitsandbytes con el esquema NF4 (\textit{Normal Float 4}),
que cuantiza los pesos a cuatro bits usando una distribución normal como grid de cuantización. Adicionalmente, se activa la doble cuantización (\textit{double quant}), que cuantiza
también las constantes de cuantización del primer nivel, reduciendo el consumo de memoria en aproximadamente 0.4 bits por parámetro. El tipo de dato de cómputo durante el paso
forward es \texttt{float16}. Esta combinación permite mantener el modelo base en memoria con un consumo inferior a cuatro gigabytes, dejando margen para los adaptadores LoRA, el
optimizador y los activaciones intermedias.

\begin{table}[htbp]
\centering
\caption{Arquitectura del modelo base Qwen2.5-7B-Instruct.}
\label{tab:AnexoA_model}
\begin{tabular}{ll}
\toprule
\textbf{Parámetro} & \textbf{Valor} \\
\midrule
Arquitectura & \texttt{Qwen2ForCausalLM} (decoder-only) \\
Número de capas & 28 \\
Tamaño del espacio oculto (\texttt{hidden\_size}) & 3584 \\
Cabezas de atención & 28 (GQA con 4 cabezas KV) \\
Tamaño del estado intermedio & 18944 \\
Función de activación & SiLU \\
Ventana de contexto máxima & 32768 tokens \\
Método de cuantización & NF4 + doble cuantización \\
Tipo de dato de cómputo & \texttt{float16} \\
\bottomrule
\end{tabular}
\end{table}

\section{Configuración LoRA}

Todas las variantes de ajuste fino (\textit{Serialized SFT}, \textit{Two-Head Model} y sus derivados) comparten la misma configuración de LoRA (\textit{Low-Rank Adaptation}).
Los adaptadores se insertan en las cuatro proyecciones de atención del transformer: \texttt{q\_proj}, \texttt{k\_proj}, \texttt{v\_proj} y \texttt{o\_proj}. No se adaptan los sesgos
ni las capas de alimentación directa (\textit{feedforward}).

\begin{table}[htbp]
\centering
\caption{Configuración LoRA común a todos los experimentos de ajuste fino.}
\label{tab:AnexoA_lora}
\begin{tabular}{ll}
\toprule
\textbf{Parámetro} & \textbf{Valor} \\
\midrule
Rango ($r$) & 8 \\
Factor de escala ($\alpha$) & 16 \\
Escala efectiva ($\alpha / r$) & 2.0 \\
Dropout & 0.05 \\
Módulos objetivo & \texttt{q\_proj}, \texttt{k\_proj}, \texttt{v\_proj}, \texttt{o\_proj} \\
Adaptación de sesgos & ninguna (\texttt{bias = none}) \\
Tipo de tarea & \texttt{CAUSAL\_LM} \\
\bottomrule
\end{tabular}
\end{table}

La elección del rango $r = 8$ responde a un compromiso entre capacidad de adaptación y carga de memoria. Con $r = 8$, cada adaptador introduce dos matrices de bajo rango de
dimensiones $3584 \times 8$ y $8 \times 3584$, lo que suma aproximadamente 230000 parámetros entrenables por capa de atención. El número total de parámetros ajustados es, en
consecuencia, pequeño en comparación con los 7000 millones del modelo base, lo que hace viable el entrenamiento en la GPU disponible. Rangos mayores aumentarían la capacidad
expresiva pero incrementarían también el consumo de VRAM durante el paso de gradientes.

\section{Hiperparámetros de entrenamiento por variante}

Este anexo recoge los hiperparámetros de los experimentos que se han ejecutado o documentado durante la evolución del trabajo. No todos ocupan el mismo lugar en la comparación
principal: algunos son ramas centrales, otros son ablaciones o ajustes diagnósticos, y otros corresponden a extensiones posteriores como DPO. La tabla~\ref{tab:AnexoA_train_common}
recoge la configuración común de los entrenamientos supervisados, mientras que las tablas posteriores separan las decisiones específicas de cada familia experimental.

\begin{table}[htbp]
\centering
\caption{Parámetros de entrenamiento comunes a las variantes SFT supervisadas.}
\label{tab:AnexoA_train_common}
\begin{tabularx}{\textwidth}{p{0.32\textwidth}X}
\toprule
\textbf{Parámetro} & \textbf{Valor} \\
\midrule
Muestras de entrenamiento & 426 (213 muestras $\times$ 2 variantes de prompt) \\
Muestras de validación & 70 (35 muestras $\times$ 2 variantes de prompt) \\
Tamaño de batch & 1 \\
Pasos de acumulación de gradiente & 8 (batch efectivo = 8) \\
Tasa de aprendizaje base & $2 \times 10^{-4}$ \\
Longitud máxima de secuencia & 2048 tokens \\
Política de tokens de entrada & tokens de prompt enmascarados con $-100$; tokens de salida supervisados \\
Orden de entrenamiento & secuencial (sin \textit{shuffle}) \\
Semilla aleatoria & 42 \\
Recuperación ante OOM & omitir el microbatch problemático cuando la variante lo permite \\
\bottomrule
\end{tabularx}
\end{table}

\begin{table}[htbp]
\centering
\small
\caption{Inventario de experimentos respecto al contenido previo del anexo.}
\label{tab:AnexoA_inventory}
\begin{tabularx}{\textwidth}{p{0.27\textwidth}p{0.23\textwidth}X}
\toprule
\textbf{Experimento} & \textbf{Estado en el anexo previo} & \textbf{Situación tras la auditoría} \\
\midrule
\texttt{baseline} zero-shot & documentado en inferencia & Se mantiene como referencia sin hiperparámetros de entrenamiento. \\
\texttt{serialized\_sft} & documentado como una época & Se corrige al entrenamiento principal de 3 épocas y se deja constancia de la ablación de una época. \\
\texttt{two\_head\_sft} & documentado & Se mantiene como rama supervisada principal con cabezal categórico. \\
\shortstack[l]{\texttt{two\_head\_ordinal}\\\texttt{\_density}} & no documentado & Se añade como primera familia ordinal con pérdida ponderada por densidad. \\
\shortstack[l]{\texttt{two\_head\_ordinal}\\\texttt{\_positional}} & documentado parcialmente & Se mantiene y se precisa como variante V1 con concatenación posicional final. \\
\shortstack[l]{\texttt{two\_head\_ordinal\_film}\\\texttt{\_positional}} & no documentado & Se añade como extensión posterior ejecutada con modulación FiLM. \\
\texttt{serialized\_sft\_dpo} & no documentado & Se añaden las ejecuciones DPO completas con \(\beta=0{,}10\) y \(\beta=0{,}30\), y se documenta la generación del dataset DPO v2. \\
\bottomrule
\end{tabularx}
\end{table}

\begin{table}[htbp]
\centering
\caption{Parámetros específicos de \texttt{serialized\_sft}.}
\label{tab:AnexoA_sft}
\begin{tabularx}{\textwidth}{p{0.34\textwidth}X}
\toprule
\textbf{Parámetro} & \textbf{Valor} \\
\midrule
Épocas de entrenamiento & 3 \\
Checkpoint principal seleccionado & \texttt{checkpoint-step-159} \\
Checkpoint de ablación de una época & \texttt{checkpoint-step-53} \\
Máximo de tokens nuevos (inferencia) & 512 \\
Formato de salida objetivo & \texttt{blocks\_tsv\_v1} \\
Cabezal estructural & ninguno (nivel codificado en la superficie textual) \\
\bottomrule
\end{tabularx}
\end{table}

\begin{table}[htbp]
\centering
\caption{Parámetros específicos de \texttt{two\_head\_sft}.}
\label{tab:AnexoA_twohead}
\begin{tabularx}{\textwidth}{p{0.34\textwidth}X}
\toprule
\textbf{Parámetro} & \textbf{Valor} \\
\midrule
Épocas de entrenamiento & 3 \\
Checkpoint seleccionado & \texttt{checkpoint-step-159} \\
Máximo de tokens nuevos (inferencia) & 1024 \\
Formato de salida objetivo & \texttt{content\_blocks\_v1} \\
Política de alineamiento de nivel & \texttt{record\_prefix\_state} \\
\midrule
\multicolumn{2}{l}{\textit{Cabezal de nivel (clasificación MLP)}} \\
Tipo & MLP clasificador \\
Número de clases & 9 (niveles 0--8) \\
Dimensión oculta & 256 \\
Dropout & 0.05 \\
Función de pérdida & entropía cruzada \\
Peso relativo de la pérdida de nivel ($\lambda_\text{level}$) & 1.0 \\
\bottomrule
\end{tabularx}
\end{table}

\begin{table}[htbp]
\centering
\small
\caption{Serie de variantes ordinales previas a la incorporación explícita de posición.}
\label{tab:AnexoA_ordinal_series}
\begin{tabularx}{\textwidth}{p{0.23\textwidth}p{0.22\textwidth}X}
\toprule
\textbf{Run} & \textbf{Diferencia principal} & \textbf{Hiperparámetros específicos} \\
\midrule
\texttt{density-v2} (2026-05-19) &
Primer cabezal ordinal con ponderación por densidad &
3 épocas; \texttt{checkpoint-step-159}; $\lambda_\text{level}=2.0$; pérdida \textit{density-weighted ordinal BCE}; 8 umbrales aprendibles; LR base $2\times10^{-4}$. \\
\texttt{thr-lr25} (2026-05-20) &
Aumento de aprendizaje de umbrales &
Misma configuración base; multiplicador de LR de umbrales $\times25$; LR de umbrales $5\times10^{-3}$. \\
\texttt{mlp3-thr50} (2026-05-22) &
Aumento conjunto del proyector ordinal y los umbrales &
Multiplicador del proyector MLP $\times3$; LR del proyector $6\times10^{-4}$; multiplicador de umbrales $\times50$; LR de umbrales $1\times10^{-2}$. \\
\texttt{gap05-mlp3-thr50} (2026-05-23) &
Inicialización centrada de los umbrales &
Misma configuración de LR diferenciadas; centro inicial 0.0; separación 0.5; 8 umbrales equidistantes entre $-1.75$ y $1.75$. \\
\bottomrule
\end{tabularx}
\end{table}

\begin{table}[htbp]
\centering
\caption{Parámetros específicos de las variantes ordinales con información posicional.}
\label{tab:AnexoA_ordinal}
\begin{tabularx}{\textwidth}{p{0.34\textwidth}X}
\toprule
\textbf{Parámetro} & \textbf{Valor} \\
\midrule
Épocas de entrenamiento & 3 \\
Máximo de tokens nuevos (inferencia) & 1024 \\
Formato de salida objetivo & \texttt{content\_blocks\_v1} \\
Política de alineamiento de nivel & \texttt{record\_prefix\_state} \\
\midrule
\multicolumn{2}{l}{\textit{Tasas de aprendizaje diferenciadas}} \\
LM base (LoRA) & $2 \times 10^{-4}$ \\
Proyector MLP ordinal & $6 \times 10^{-4}$ (multiplicador $\times 3$) \\
Umbrales aprendibles & $1 \times 10^{-2}$ (multiplicador $\times 50$) \\
\midrule
\multicolumn{2}{l}{\textit{Cabezal ordinal con codificación posicional}} \\
Tipo V1 & \texttt{learned\_threshold\_ordinal\_positional} \\
Tipo V2 & \shortstack[l]{\texttt{learned\_threshold\_ordinal}\\\texttt{\_film\_positional}} \\
Número de clases & 9 (niveles 0--8) \\
Dimensiones del proyector & [512, 64] \\
Dropout por capa & [0.1, 0.0] \\
Codificación posicional & sinusoidal absoluta, $d=16$, 8 frecuencias $\{1, 2, 4, 8, \ldots, 128\}$ \\
Inyección posicional V1 & concatenación final (\textit{final\_concat}, causal) \\
Inyección posicional V2 & modulación FiLM tras la capa de 512 dimensiones (\textit{film\_after\_512}) \\
Configuración FiLM V2 & dimensión oculta 128; escala inicial 0.1; última capa inicializada a identidad \\
Función de pérdida & \textit{density-weighted ordinal BCE} \\
Peso relativo de la pérdida de nivel ($\lambda_\text{level}$) & 2.0 \\
Número de umbrales aprendibles & 8 ($\tau_1, \ldots, \tau_8$ para 9 clases) \\
Parametrización de umbrales & $\tau_0 + \sum_{i} \text{softplus}(\delta_i)$ \\
Inicialización de umbrales & equidistante, centro 0, separación 0.5 \\
\bottomrule
\end{tabularx}
\end{table}

\begin{table}[htbp]
\centering
\caption{Parámetros específicos de las ejecuciones DPO sobre \texttt{serialized\_sft}.}
\label{tab:AnexoA_dpo}
\begin{tabularx}{\textwidth}{p{0.34\textwidth}X}
\toprule
\textbf{Parámetro} & \textbf{Valor} \\
\midrule
Variante de partida & \texttt{serialized\_sft}, \texttt{checkpoint-step-159} \\
Formato de preferencia & pares \texttt{chosen}/\texttt{rejected} en \texttt{blocks\_tsv\_v1} \\
Muestras de preferencia & 454 \\
Candidatos por prompt para preferencias & 6 \\
Temperaturas de muestreo de candidatos & 0.45, 0.65, 0.85, 1.0, 1.1, 1.2 \\
\texttt{top\_p} en generación de candidatos & 0.95 \\
Máximo de tokens nuevos en generación de candidatos & 512 \\
Batch y semilla en generación de candidatos & batch 1; semilla 42 \\
Épocas de entrenamiento & 1 \\
Tamaño de batch & 1 \\
Pasos de acumulación de gradiente & 8 \\
Tasa de aprendizaje & $5\times10^{-6}$ \\
Longitud máxima de secuencia & 2048 tokens \\
Máximo de tokens nuevos & 1024 \\
Pasos entre checkpoints & 32 \\
Warmup & 0.03 \\
Decaimiento de pesos & 0.0 \\
Gradient checkpointing & activado \\
Recuperación ante OOM & fallar la ejecución \\
Configuraciones \(\beta\) ejecutadas & \(0{,}10\) completa; \(0{,}30\) completa \\
Política de referencia & modelo \texttt{serialized\_sft} congelado; log-probabilidades de referencia precomputadas \\
\bottomrule
\end{tabularx}
\end{table}

\begin{table}[htbp]
\centering
\caption{Parámetros de generación y selección del dataset de preferencias DPO v2.}
\label{tab:AnexoA_dpo_v2_dataset}
\begin{tabularx}{\textwidth}{p{0.34\textwidth}X}
\toprule
\textbf{Parámetro} & \textbf{Valor} \\
\midrule
Política generadora & checkpoint final de DPO v1 con \(\beta=0{,}30\) (\texttt{checkpoint-step-57}) \\
Política de referencia prevista para la siguiente iteración & misma política generadora, congelada como \(\pi_{\mathrm{ref}}^{(2)}\) \\
Split usado para generar candidatos & entrenamiento de \texttt{kubernetes\_v1} \\
Prompts considerados & 426 \\
Candidatos por prompt & 6 \\
Candidatos evaluados & 2556 \\
Temperaturas de muestreo de candidatos & 0.45, 0.65, 0.85, 1.0, 1.1, 1.2 \\
\texttt{top\_p} en generación de candidatos & 0.95 \\
Máximo de tokens nuevos en generación de candidatos & 512 \\
Batch y semilla en generación de candidatos & batch 1; semilla 42 \\
Pares finales seleccionados & 229 \\
Prompts con al menos un par & 193 \\
Tipos de pares & \texttt{domain\_invariant}, \texttt{gate\_crossing}, \texttt{level5\_practice}, \texttt{prompt\_fidelity}, \texttt{structural\_fidelity} \\
Margen mínimo de preferencia & 0.25 \\
Restricciones de conservación & no degradar en exceso fidelidad al prompt, F1 de líneas, exactitud de \texttt{level} ni campos obligatorios \\
Estado del dataset & generado, analizado y usado en dos entrenamientos DPO v2 con evaluación de validación completa o parcialmente evaluable según parseo estructurado \\
\bottomrule
\end{tabularx}
\end{table}

\begin{table}[htbp]
\centering
\caption{Parámetros específicos de las ejecuciones DPO v2 sobre la política DPO v1.}
\label{tab:AnexoA_dpo_v2_training}
\begin{tabularx}{\textwidth}{p{0.34\textwidth}X}
\toprule
\textbf{Parámetro} & \textbf{Valor} \\
\midrule
Política inicial y de referencia & DPO v1 con \(\beta=0{,}30\), \texttt{checkpoint-step-57} \\
Dataset de preferencias & \texttt{agent-full-auto-v2/preferences\_final.jsonl} \\
Pares de preferencia & 229 \\
Formato de salida & \texttt{blocks\_tsv\_v1} \\
Épocas de entrenamiento & 3 \\
Tamaño de batch & 1 \\
Pasos de acumulación de gradiente & 8 \\
Configuración \(\beta\) & \(0{,}30\) \\
Tasas de aprendizaje ejecutadas & \(1\times10^{-5}\) y \(1\times10^{-4}\) \\
Longitud máxima de secuencia & 2048 tokens \\
Máximo de tokens nuevos en validación & 1024 \\
Pasos entre checkpoints & 29 \\
Checkpoint final evaluado & \texttt{checkpoint-step-87} en ambas variantes \\
Política de referencia & log-probabilidades precomputadas antes de cargar la política entrenable \\
\bottomrule
\end{tabularx}
\end{table}

\section{Configuración de inferencia y evaluación}

Durante la evaluación, todos los modelos generan texto en modo greedy, es decir, seleccionando en cada paso el token con mayor probabilidad sin muestreo. Esta elección
elimina la varianza introducida por la temperatura y hace que los resultados sean deterministas, lo que facilita la comparación directa entre variantes. El límite de tokens
nuevos varía según el formato de salida esperado: la línea base y el SFT serializado generan representaciones más compactas, mientras que el two-head genera únicamente contenido
y metadatos de posición sin incluir el nivel en la cadena de texto, lo que puede requerir secuencias más largas para completar el mismo número de líneas.

Esta configuración greedy se refiere a la evaluación comparativa de los modelos ya entrenados. La generación de candidatos para DPO se trata de forma separada: ahí sí se utiliza
muestreo con varias temperaturas para obtener respuestas suficientemente distintas sobre un mismo prompt. Esta diferencia no introduce una segunda política de evaluación, sino que
responde a dos usos distintos de la generación: medir modelos de forma determinista, por un lado, y construir contrastes de preferencia informativos, por otro.

\begin{table}[htbp]
\centering
\caption{Configuración de inferencia por variante experimental.}
\label{tab:AnexoA_inference}
\begin{tabular}{lccc}
\toprule
\textbf{Parámetro} & \textbf{Baseline} & \textbf{serialized\_sft} & \textbf{two\_head (todas)} \\
\midrule
Temperatura & 0.0 & 0.0 & 0.0 \\
top\_p & 1.0 & 1.0 & 1.0 \\
Máximo de tokens nuevos & 320 & 512 & 1024 \\
Batch de inferencia & 1 & 1 & 1 \\
\bottomrule
\end{tabular}
\end{table}

\section{Reproducibilidad y trazabilidad}

El diseño del sistema de entrenamiento incorpora varios mecanismos para garantizar que los experimentos puedan reproducirse o reanudarse sin pérdida de estado.
En primer lugar, la semilla aleatoria se fija a 42 en todos los experimentos, tanto para la inicialización de los adaptadores LoRA como para los procesos estocásticos
internos de PyTorch y Python. En segundo lugar, el conjunto de entrenamiento se procesa en orden secuencial sin mezcla, de modo que el orden de presentación de los ejemplos
es idéntico en cada ejecución. En tercer lugar, los puntos de control almacenan no solo los pesos del adaptador sino también el estado completo del optimizador y el planificador
de tasa de aprendizaje, lo que permite reanudar el entrenamiento exactamente desde el último paso guardado.

Cuando un microbatch provoca un error de memoria de GPU (\textit{CUDA out of memory}), el sistema lo registra en el fichero \texttt{oom\_skipped\_batches.jsonl}, lo excluye
del paso de optimización y continúa con el siguiente. Este mecanismo evita que un único ejemplo problemático interrumpa la ejecución, a costa de que ese ejemplo no contribuya
a la actualización del modelo en ese paso. Finalmente, todas las ejecuciones de entrenamiento y validación intermedia se registran en la plataforma Weights \& Biases bajo el
proyecto \texttt{llm-structured-semantic-generation}, con logs de pérdida, métricas de validación por paso de \textit{checkpoint} y configuración completa del experimento.

\section{Ejemplos de entrada y salida del mejor sistema}

Para complementar la comparación agregada del Capítulo~\ref{ch::chapter8}, esta sección recoge tres ejemplos concretos generados por la variante DPO v2 evaluada sobre el split de test. En todos los casos se usa el \textit{checkpoint} \texttt{checkpoint-step-87}; la inferencia se realiza en modo greedy y el modelo produce internamente una secuencia \texttt{blocks\_tsv\_v1}. Por tanto, las salidas que se muestran no son una edición manual posterior, sino el YAML reconstruido por el parser a partir de esos bloques estructurados.

Los ejemplos se han seleccionado para mostrar tres tipos de petición distintos: un objeto RBAC compacto, una modificación estructural sobre un \texttt{DaemonSet} y un recurso \texttt{Endpoints} relacionado con la comunicación entre espacios de nombres. Su objetivo no es sustituir a las métricas, sino hacer visible cómo se materializa la salida final del sistema en casos concretos del conjunto de test.

\subsection*{Ejemplo 1: RoleBinding RBAC compacto}

\noindent\textit{Identificador:} \texttt{q224::question\_simplified}.

\noindent\textbf{Prompt.}
\begin{lstlisting}[style=yamlblock]
Craft a k8s yaml for RoleBinding "testadminbinding" in "test3" ns. Grant "myaccount" SA from "test" ns "cluster-admin" CR permissions.
\end{lstlisting}

\noindent\textbf{Predicción.}
\begin{lstlisting}[style=yamlblock]
apiVersion: rbac.authorization.k8s.io/v1
kind: RoleBinding
metadata:
  name: testadminbinding
  namespace: test3
roleRef:
  apiGroup: rbac.authorization.k8s.io
  kind: ClusterRole
  name: cluster-admin
subjects:
- apiGroup: ''
  kind: ServiceAccount
  name: myaccount
  namespace: test
\end{lstlisting}

\subsection*{Ejemplo 2: DaemonSet con \texttt{nodeSelector}}

\noindent\textit{Identificador:} \texttt{q18::question\_simplified}.

\noindent\textbf{Prompt.}
\begin{lstlisting}[style=yamlblock]
Add a nodeSelector to the given DaemonSet YAML to ensure Pods run only on nodes labeled "disktype=ssd".
\end{lstlisting}

\noindent\textbf{Predicción.}
\begin{lstlisting}[style=yamlblock]
apiVersion: apps/v1
kind: DaemonSet
metadata:
  name: nodeselector-daemonset
spec:
  selector:
    matchLabels:
      app: nodeselector-daemonset
  template:
    metadata:
      labels:
        app: nodeselector-daemonset
    spec:
      containers:
      - command:
        - sleep
        - '3600'
        image: busybox
        name: busybox
      nodeSelector:
        disktype: ssd
\end{lstlisting}

\subsection*{Ejemplo 3: Endpoints entre espacios de nombres}

\noindent\textit{Identificador:} \texttt{q257::question\_simplified}.

\noindent\textbf{Prompt.}
\begin{lstlisting}[style=yamlblock]
How can I define a service in one namespace that links to a Pod running in another namespace without the code inside the container needing to know about the location of the service?
\end{lstlisting}

\noindent\textbf{Predicción.}
\begin{lstlisting}[style=yamlblock]
apiVersion: v1
kind: Endpoints
metadata:
  name: service-x
  namespace: namespace-a
subsets:
- addresses:
  - targetRef:
      kind: Service
      namespace: namespace-b
      name: service-x
      apiVersion: v1
  ports:
  - port: 80
\end{lstlisting}
